\section{Paths Forward} \label{sec:main-paths}

Now, we describe several paths forward to address the aforementioned challenges. We group these into three parts: data collection, training, and inference time approaches. 

\subsection{Data Collection} 

When developing LLMs for code, the open-source community relies on datasets like \textit{the Stack} \citep{lozhkov2024starcoder}, consisting of trillions of GitHub tokens. However, richer sources of data exist such as fine-grained data of the developmental process. As this data needs to be collected on a large scale, we call for a community-based open source effort to assemble an open dataset for training code LLMs.

\textbf{Automatic Data Curation}: The advantage of code is it is possible to achieve strong, verifiable feedback with test cases, program execution engines, and other symbolic tools. Modern programming tools allow us to extract rich semantic and structural information about code. 

\textit{Augmenting Data with Program Information}: One way to increase semantic richness of coding datasets is to augment training datasets with detailed annotations describing various properties of programs. We hypothesize that this augmentation will significantly improve a model’s understanding of code, leading to better generalization and stronger coding capabilities. This information includes static analysis (abstract syntax trees), program instrumentation, dynamic analysis, and formal verification. 

\textit{High-quality, Verifiable Synthetic Data}: The verifiability of code makes it possible to generate a large batch of data and filter out low-quality samples, even at the repository level. With the supervision of experts in various domains, we believe the community can come up with high-quality synthetic data pipelines with broad coverage of software engineering tasks. For example, to generate code with interesting program invariants, we can sample a large batch of programs, run an invariant detection engine, and retain only programs with interesting invariants.

\textbf{Human-Centric Data Curation}: While automatic data curation can provide a lot of additional information, there are still signals that are best obtained via manual efforts. We list three such categories of tasks below:

\textit{Fine-Grained Data of the Developmental Process}: IDE providers have an advantage when it comes to coding data, as they have direct access to the full history and evolution of a codebase. Companies such as Google and Meta \citep{chandra2024ai, murali2024ai} collect fine-grained signals such as edit history. With direct access to the full history and evolution of a codebase, they can track which suggestions are adopted over time. However, companies generally collect data for their own use cases, making an open-source effort invaluable for curating a collective and open use dataset.

\textit{Data for Diverse SWE Tasks}: Most of today's code LLM training recipes focus on code generation because large-scale datasets are mostly in a continuous format. We believe that involving domain experts and curating data for all sorts of tasks (including challenging and out-of-distribution ones) will lead to diversity and more general coding capabilities.

\textit{Human-Centric Data}: Code LLMs are trained and evaluated on carefully curated datasets with clear instructions and verifiable test cases. However, these models are often deployed in real-world scenarios where users provide vague specifications or incomplete requirements in their queries. Collecting human-centric data reflecting real-world model usage is a promising approach to bridging the gap between model development and deployment, but we currently have very little data of this form. An open-source community effort can lead to tools and environments such as open-source IDEs to capture real-world interactions in diverse modalities.

\subsection{Training}

\textbf{Environment Design for Code RL}: In recent months, RLVR has seen success in solving algorithmic programming problems through DeepSeek-R1~\citep{deepseekai2025deepseekr1} and OpenAI o1. Recently, on SWE-Bench, SWE-RL \citep{wei2025swerladvancingllmreasoning} use RL on a rule-based reward to improve performance on SWE-Bench. We find it promising to continue scaling the RL approach and create environments from tasks collected from real-world software engineering repositories. These environments can be used to improve reasoning skills, environment-interaction capabilities and tool usage. However, scaling this up significantly requires solving several research and engineering problems. 
First, installing arbitrary repositories from Github, even using CI is challenging and we require smarter solutions potentially involving LLM-based installation agents.
Next, setting up execution infrastructure would require storing installed repository images in something akin to \verb|docker| for efficient storage and fast container startup times~\citep{team2025kimi}.
Notably, combined docker images can grow massively large and often grow at hundreds of gigabytes even at a modest scale of a few hundred repositories. Because of the scale and engineering effort required, we advocate for the community to come together and create a unified gym environment for coding.

% Beyond environments, performing large-scale reinforcement learning would require collecting diverse challenging problems with an appropriate way to compute the rewards. 
% These task prompts can be collected from Github~\citep{pan2024trainingsoftwareengineeringagents} or generated synthetically from problems on Github.
% Moreover, assuming access to many executable repositories, we can source various end-to-end problems for tasks beyond bug-fixing such as optimization, fuzzing, etc.
% Access to pre-existing or generated test cases allows for measuring correctness and providing rewards. 

% However, we envision many practical challenges to remain. 
% % \textbf{Task ambiguity.}
% For example, longer-horizon tasks are usually more ambiguous and approaches may require multi-turn interactions beyond autonomous coding agents.
% % most long-horizon tasks will be ambiguous and will require multi-turn interactions.
% This would pose a considerable challenge during reinforcement learning where ambiguity resolution might need to be modeled in the reinforcement learning process itself. 
% We elaborate on human collaboration further in Section~\ref{sec:d-hai-training}.
% % \textbf{Reward Computation.}
% % Testing is key for reward computation. We can generate tests (similar to r2e). 
% % Reward Hacking is still a concern. 
% % Cite negative examples, some from r2e, some from sakana
% Reward hacking~\cite{skalse2022defining} poses another challenge as we build more real-world coding challenges. 
% Test cases often suffer from coverage issues and can grade correct solutions as incorrect.
% For example, ~\cite{baker2025monitoring,denison2024sycophancy} identified that models attempt to bypass or cheat against the testing harness when optimized using reinforcement learning. 

% \textbf{Rewards without execution}: As setting up execution environments can lead to considerable overhead, another potential strategy is to use proxy metrics and trained language models to judge correctness. This was common in the pre-LLM era, researchers often used BLEU/CodeBLEU \citep{papineni2002bleu, ren2020codebleu} and BERTScore/CodeBERTScore \citep{zhang2019bertscore, zhou2023codebertscore} to assess correctness of text and code. In code, semantic and structural properties can be used to improve similarity metrics. Two examples of this are Dolos \citep{maertens2022dolos}, an AST-aware plagiarism detector, and \texttt{difflib.SequenceMatcher}, which can be used to compute the similarity between two patches \citep{wei2025swerladvancingllmreasoning, ma2025sorft}. 
% Beyond rule-based rewards, LLMs-as-a-judge approaches can also be used as reward functions, possibly in conjunction with other execution-based or execution-free approaches.



\textbf{Adapting to Specialized and Quickly Changing Codebases}: Being able to quickly adapt to code is an important and crucial skill, as codebases are often out-of-distribution and frequently change. Here, we describe three potential research directions leading to such adaptation. 

\textit{Test-time training (TTT) to custom codebases}: TTT is the recent paradigm of adapting to a specific problem instance by training on a narrow set of in-distribution examples~\citep{akyurek2024surprisingeffectivenesstesttimetraining,sun19ttt}. This can be used when working in a low-resource context, for example training on a specific codebase, new domain, or unseen API. One way to get data in these contexts is to collect and learn from trajectories of a SWE agent's behavior. We can keep track of previous model attempts and failures to prevent repeated mistakes. This will steer the model closer to the desired distribution, such as generating code in the specified version of libraries in the current context.

\textit{Keeping an information bank of code information}: For library and versioning issues, retrieval can be very effective for preventing hallucinations of wrong versions of libraries, which can inherently lead to better synthetic data and agentic trajectories. During the TTT process, we can also keep a large growing memory bank of code, documentation, synthetic code, and agentic trajectories in the specialized context. Retrieving from the memory bank would improve the success of generating code, which can then be augmented to the memory bank, and so on, continuously increasing the amount of data and knowledge.

\textit{Learning on the fly}: When humans are faced with a task they have never seen before, they are often able to draw from past experiences and quickly adapt and generalize to the new domain. This is one of the big unsolved challenges of today's LLMs: given an OOD coding task, how can models get up to speed and productively work on the task with few samples? On toy domains, an example of this is DreamCoder \citep{ellis2021dreamcoder}, a system that learns to solve problems by writing programs and automatically discovering domain concepts and composing concepts together. Designing such approaches for more practical applications is an exciting research direction that will have drastic implications for coding and reasoning.

\textbf{Training Code LLMs to Collaborate with Humans}: As we collaborate with AI more to write code, it is crucial that these models have capabilities that allow for effective collaboration. 

\textit{Specifications Beyond Natural Language}: While natural language prompts offer intuitive and flexible ways to express requirements, they are often ambiguous and incomplete. To address this, we can train models to use more precise and verifiable representations such as formal specifications and test-based specifications. Formal specifications allow an AI system to precisely verify correctness. Test-based specifications range from input-output examples and assertions to property-based tests. However, hand-crafted test suites can be incomplete, leading to misalignment where AI-generated code passes tests but does not genuinely meet functional requirements. Moving forward, we should train models to generate high-quality test cases based on the user's initial query, ensuring more comprehensive specification coverage.

\textit{Learning to Quantify Uncertainty and Communicate Proactively}: As AI coding systems are increasingly deployed to complex SWE tasks, they encounter more ambiguous and uncertain scenarios compared to traditional benchmarks for coding models. Ideally, in such situations, these systems should proactively communicate with users to clarify tasks and acknowledge its own limitations rather than becoming stuck in endless failure loops or generating buggy code. A key challenge is enabling models to distinguish between well-specified and ambiguous instructions while quantifying uncertainty in a robust manner, which will require incorporating corresponding reasoning data into post-training. Another challenge in human-agent collaboration is communication \citep{shao2024collaborativegym}, highlighting the need to improve the proactive communication capability of the models. Current models often fail to ask meaningful questions when user input is ambiguous or insufficient and struggle to provide progress updates or verify plans in interactive settings. Enhancing the proactive communication skills of models requires innovative approaches to reward behaviors that yield benefits over multiple steps. Since communication with users does not immediately resolve the task at hand, but may improve long-term results, effective strategies must account for delayed rewards in training.

\subsection{Inference Time Approaches}

In this section, we advocate for several inference-time research thrusts that have currently been underexplored. Making progress on these fundamental directions can unblock many of the challenges previously mentioned.

\textbf{Semantic-Aware Embeddings and Retrieval}: We saw that code that is close in embedding space is more often syntactically similar than semantically similar \citep{zhao2023get}, making it hard to retrieve semantically similar code. However, before the LLM era, there were a variety of efforts to incorporate code properties when training embeddings. For example, \citet{nye2020representing} train neural modules to represent program operations, leading to compositional program representations that encode the semantics of the underlying programming language. Many other works \citep{zohar2018automatic, ellis2019write, chen2021latent} attempt to learn execution-aware latent representations for partial and full programs, taking semantics into account. 

Incorporating these techniques to train models to have better and more semantically aware representations can lead to models with a more general understanding of code (Sec. \ref{sec:c-global-understanding}). For example, if correct and buggy programs could hypothetically be separated in embedding space, then models could be steered away from the incorrect program space. While a complete separation might not be possible, having semantics-aware embeddings could lead to downstream improvements on SWE tasks.

\textit{Retrieving via on-the-fly navigation}: Instead of keeping track of embeddings, another approach is to find retrievals on-the-fly by navigating the codebase. This prevents the high cost of continuously maintaining and updating the retrieval bank. An agent can learn to use command line functions such as \texttt{cd}, \texttt{ls}, and \texttt{grep}, and IDE functions such as jumping to function definitions or finding all references of a function. Static analysis tools can also be paired with the agent to improve code navigation, such as providing the abstract syntax tree (AST) or file structure of a codebase.

\textbf{Integration with SWE Development Frameworks}: As AI improves at SWE, there are increasingly more opportunities to incorporate it into the continuous integration and continuous deployment (CI/CD) process. In CI/CD, automated pipelines are the backbone for building, testing, and deploying code changes. CI/CD accelerates feedback cycles and minimizes integration issues. AI could contribute in many ways. First, AI-powered code review tools can be incorporated into CI pipelines to automatically identify and flag style violations, potential security vulnerabilities, and code smells before human reviewers are involved. Second, AI can provide intelligent deployment risk assessments. By analyzing code changes, test outcomes, and historical deployment data, AI can predict the likelihood of deployment issues, informing decisions about whether to proceed with automated deployment or mandate manual verification steps. Finally, AI can automate the generation of release notes by summarizing commit messages, issue tracker data, and relevant code modifications within the CI/CD process.

\textit{Steering away from software anti-patterns}: In SWE, certain anti-patterns frequently lead to bugs. Common weakness enumeration (CWE) is a categorization of software and hardware weaknesses often leading to vulnerabilities. Because publicly available GitHub code often contains code with anti-patterns, bugs, and CWEs, LLMs often write code susceptible to these issues \citep{asare2023github, fu2023security}. Explicitly steering models against these vulnerabilities can lead to more secure and correct code, e.g. by collecting programs violating each CWE (synthetically or via GitHub) and using them as negative signal during SFT or RL.

\textbf{Neurosymbolic Approaches}: Code is a unique domain because there is a vast body of techniques from programming languages (PL) research to build off of, but the majority of AI for code research today does not leverage the symbolic properties of code. Traditional PL approaches have a few common shortcomings. First, they often require very complete and precise specifications. Many tools need to have specifications for all library functions, need to specialize to a precise version of the language, and need to specialize to the build system. Second, there is often a high computational cost due to the large search space. Third, there can be many false positives due to the limitations of the tool. We believe that deeply integrating these symbolic tools with LLMs can partially mitigate these challenges. 

We provide a few examples of this potential integration. When generating code, program analysis techniques could be applied on shorter snippets of AI-generated code to surface potential bugs or prove properties. To improve general code understanding, LLMs can be trained with information about program structure such as ASTs \citep{gong2024ast}. When debugging a large codebase, AI could be first used to narrow down potentially problematic sections of the code which are then handed off to PL tools for debugging. During code generation in DSLs, LLMs can leverage the grammar of the programming language to do constrained decoding~\citep{poesia2022synchromesh, geng2023grammar, wei2023copiloting} to mitigate syntactic errors.

\textbf{Scaffolding Human Supervision}: Once code LLMs are deployed for inference, it is crucial to scaffold human supervision of AI-generated code. This goes beyond merely enhancing the accuracy of AI-generated code, as humans often still need to make the final decision on whether to accept the code or understand it for future integration and maintenance. A study on Github Copilot usage~\citep{10.1145/3551349.3560438} revealed that programmers tend to allocate less visual attention to AI-generated code. %, which further highlights the importance of designing AI systems that scaffold human oversight while reducing cognitive load during code review.
While we could train humans to better identify issues in AI-generated code~\citep{10.1145/3627217.3627238}, a more desirable approach is to design AI systems that scaffold human supervision, reducing human cognitive load when reviewing generated code.

One way to achieve this is by enriching AI-generated content with additional contextual information. In software engineering specifically, \citet{sun2024source} highlighted the benefits of high-quality source code summarization in aiding software developers in understanding and maintaining machine-generated code. Second, interactive approaches can also enhance supervision such as \textit{Live Programming} \citep{ferdowsi2024validating}, a continuous display of the runtime values of a program. Finally, improving the readability and interpretability of AI-generated code itself presents a promising direction. Expanding on these ideas, future research should prioritize human interpretability in the design and optimization of AI coding systems, fostering greater trust and control in AI-assisted software development.
